\frontsection{Первое знакомство с машинным обучением}
\label{sec:first-steps}
\subsection*{От примеров к правилу}
Предположим, в таблице записаны площадь, число комнат и известная цена нескольких домов. Мы хотим оценить цену другого дома. В машинном обучении правило прогноза настраивают по примерам с известным ответом. Эти примеры называются \textbf{обучающими данными}, а настраиваемое правило --- \textbf{моделью}.

\begin{table}[H]\centering\small
\begin{tabularx}{\textwidth}{@{}p{0.22\textwidth}Y@{}}
\toprule
Понятие & Пример и смысл \\
\midrule
Объект & Один дом; в таблице ему соответствует одна строка. \\
Признак & Известная характеристика дома, например площадь. \\
Целевой ответ & То, что требуется предсказать: здесь цена. \\
Выборка & Набор объектов, выделенный для определённой цели. \\
Прогноз & Ответ модели для конкретного объекта. \\
Обучение & Настройка модели по признакам и известным ответам. \\
\bottomrule
\end{tabularx}
\end{table}

При \textbf{регрессии} ответ --- число, например цена. При \textbf{классификации} ответ --- одна из заранее заданных категорий, например «ходит» или «сидит». Категории называют классами, а правильный класс объекта --- его меткой. Числовой код класса служит обозначением и не превращает задачу в регрессию.

Во всех трёх работах рассматривается \textbf{обучение с учителем}: для обучающих примеров известны правильные ответы. «Учитель» здесь означает наличие ответов в данных. Для нового объекта модель получает только те сведения, которые доступны при её применении.

\subsection*{Как читать обозначения}
Запись \(x_i\) обозначает признаки объекта с номером \(i\), \(y_i\) --- правильный ответ, \(\widehat y_i\) --- прогноз. Например, \(x_i=(80,3)\) может описывать площадь и число комнат. Символ \(f_\theta\) обозначает модель, а \(\theta\) --- её настраиваемые параметры:
\[
\widehat y_i=f_\theta(x_i).
\]
Эта формула читается так: «подать признаки объекта в модель и получить прогноз». В работе с изображениями признаки на входе --- значения пикселей; в работе с активностью --- уже вычисленные характеристики сигналов.

\clearpage
\subsection*{Почему нельзя проверять только на знакомых примерах}
Правило может хорошо отвечать на знакомые примеры и ошибаться на новых. Способность работать на новых объектах называют \textbf{обобщением}. Если модель слишком подстроилась под особенности обучения и хуже переносится на новые данные, говорят о \textbf{переобучении}.

Для организации проверки данные разделяют по ролям:
\begin{enumerate}
\item \textbf{Обучение (train):} настраиваем модель.
\item \textbf{Валидация (validation):} сравниваем варианты и выбираем настройки.
\item \textbf{Итоговый тест (test):} оцениваем выбранное решение, если есть отдельные данные с ответами, не использованные при выборе.
\end{enumerate}
Можно представить подготовку к экзамену: разобранные задачи помогают учиться, пробный вариант --- выбирать способ подготовки, новый экзаменационный вариант --- проверять результат. Это аналогия ролей, а не требование делить любой набор ровно на три файла.

\textbf{Утечка данных} возникает, если при обучении используется информация, которая по условиям проверки должна быть недоступна. Например, цена дома случайно оставлена среди признаков или средние для подготовки признаков вычислены с использованием валидации.

\subsection*{Что именно меняет экспериментатор}
\textbf{Параметры} модель находит при обучении: например, веса нейросети. \textbf{Гиперпараметры} задаёт исследователь: например, число соседей или ширину скрытого слоя. \textbf{Метрика} --- число, которым оценивают качество прогноза. \textbf{Исходная модель (baseline)} --- понятный вариант, с которым сравнивают изменения.

Под \textbf{протоколом} понимается заранее записанный порядок опыта: какие данные и разбиение используются, что меняется, какая метрика сравнивается. \textbf{Бюджет} ограничивает число запусков, эпох или время. \textbf{Seed} --- начальное состояние генератора случайных чисел: оно помогает повторить выбор разбиения и начальных весов, но не делает оценку качества безошибочной.

\begin{keybox}{Проверьте себя перед первой работой}
Выберите один дом или одно изображение. Назовите объект, доступные признаки и искомый ответ. Объясните, чем настройка модели отличается от проверки её прогноза. Если ответ пока неясен, вернитесь к этим определениям до запуска кода.
\end{keybox}
